\section{Conclusion}
\label{sec:conclusion}

We tested whether knowledge editing can make an isolated change to arithmetic knowledge in \gptxl, finding that it cannot.
All three methods---\rome, fine-tuning, and \lora---successfully teach the model that $2{+}2{=}5$, but every method introduces a systematic \fivebias that contaminates arithmetic broadly.
\rome uniquely preserves general language modeling (perplexity unchanged) while still corrupting arithmetic, revealing that arithmetic relies on shared computational circuits distinct from general language circuits.

The key takeaway is that ``surgical'' knowledge edits can have predictable but hard-to-prevent side effects on computationally related knowledge.
Current locality benchmarks, which test only unrelated facts, miss these domain-specific failures.

Future work should test models with stronger arithmetic baselines, evaluate methods with theoretical preservation guarantees (\eg \alphaedit), and develop editing benchmarks that include computational knowledge alongside factual knowledge.
More broadly, mechanistic interpretability tools could trace the exact circuits responsible for arithmetic, potentially enabling edits that respect circuit boundaries.
